\section*{Chapitre \ref{ch:b3:quantum-model-systems} --- La mécanique quantique du chimiste~: systèmes modèles}

\begin{solution}{exo:b3:quantum-model-systems:1}
$E_1 = h^2/8m_eL^2$. (a) $L = \qty{1.00}{nm}$~: $E_1 = \qty{6.02e-20}{J} =
\qty{0.376}{eV}$~; $\Delta E_{12} = 3E_1$ et $\lambda = hc/3E_1 =
\qty{1099}{nm}$, dans le proche infrarouge. (b) $L = \qty{0.50}{nm}$~:
$E_1$ est quatre fois plus grande, \qty{2.41e-19}{J} $= \qty{1.50}{eV}$, et $\lambda =
\qty{275}{nm}$, dans l'ultraviolet.
\end{solution}

\begin{solution}{exo:b3:quantum-model-systems:2}
$n_x^2 + n_y^2 + n_z^2 = 3$ (1,1,1)~: $g = 1$~; $6$ (2,1,1 et permutations)~:
$g = 3$~; $9$ (2,2,1)~: $g = 3$~; $11$ (3,1,1)~: $g = 3$~; $12$ (2,2,2)~: $g = 1$.
\end{solution}

\begin{solution}{exo:b3:quantum-model-systems:3}
$[\hat x^2,\hat p]f = -\iu\hbar x^2f' + \iu\hbar(x^2f)' = 2\iu\hbar xf$, donc
$[\hat x^2,\hat p] = 2\iu\hbar\hat x$. $[\hat p^2,\hat x]f = -\hbar^2[(xf)'' -
xf''] = -2\hbar^2f'$, donc $[\hat p^2,\hat x] = -2\iu\hbar\hat p$.
\end{solution}

\begin{solution}{exo:b3:quantum-model-systems:4}
$E_J/hc = BJ(J+1)$~: 0, 21.19, 63.56 et \qty{127.12}{cm^{-1}}, avec $g = 1$,
3, 5, 7. $I = h/(8\pi^2cB) = \qty{2.643e-47}{kg.m^2}$ (avec $c$ en
\unit{cm/s}).
\end{solution}

\begin{solution}{exo:b3:quantum-model-systems:5}
$|\psi_n|^2$ est symétrique par rapport à $L/2$, donc $\langle x\rangle = L/2$.
Avec $\sin^2u = \frac12(1 - \cos2u)$ et deux intégrations par parties,
$\langle x^2\rangle = L^2\big(\frac13 - \frac{1}{2n^2\pi^2}\big)$. Pour $n = 1$,
$\langle x^2\rangle = 0.283L^2$, moins que la valeur classique $L^2/3$~:
l'état fondamental se concentre au milieu. Quand $n$ croît, on retrouve la
valeur classique.
\end{solution}

\begin{solution}{exo:b3:quantum-model-systems:6}
$1\,\unit{cm^{-1}} = \qty{0.011963}{kJ/mol}$. EPZ(\ce{H2}) $= 2200.6 \times
0.011963 = \qty{26.33}{kJ/mol}$, EPZ(\ce{D2}) $= \qty{18.63}{kJ/mol}$~;
différence
\qty{7.69}{kJ/mol}. Même \omterm{def:b3:quantum-model-systems:oscillator}{constante de force}, donc $\tilde\omega \propto \mu^{-1/2}$
et le rapport attendu est $\sqrt{\mu_D/\mu_H} = 1.4137$ (masses atomiques
2.01410 et 1.00783)~; le rapport mesuré vaut 1.4127, à 0.07\,\% près (le
petit écart vient des électrons, qui ne suivent pas parfaitement les
noyaux).
\end{solution}

\begin{solution}{exo:b3:quantum-model-systems:7}
$\int_0^Lx^2(L - x)^2\,\dd x = L^5/30$, donc $N = \sqrt{30/L^5}$. Alors
\[
\langle\psi_1|\phi\rangle = \sqrt{\tfrac2L}\sqrt{\tfrac{30}{L^5}}\int_0^Lx(L -
x)\sin\tfrac{\pi x}{L}\,\dd x = \frac{\sqrt{60}}{L^3}\cdot\frac{4L^3}{\pi^3} =
\frac{4\sqrt{60}}{\pi^3} = 0.99928 .
\]
La parabole est à 99.86\,\% l'état fondamental (le carré du recouvrement).
\end{solution}

\begin{solution}{exo:b3:quantum-model-systems:8}
$\langle f|\hat pg\rangle = -\iu\hbar\int f^*g'\,\dd x = -\iu\hbar[f^*g] +
\iu\hbar\int f^{*\prime}g\,\dd x = \int(-\iu\hbar f')^*g\,\dd x = \langle\hat
pf|g\rangle$, le crochet s'annulant aux extrémités. Pour $\hat x\hat p$, en
utilisant que $\hat x$ et $\hat p$ sont chacun \omterm{def:b3:quantum-model-systems:hermitian}{hermitiens}~:
$\langle f|\hat x\hat pg\rangle = \langle\hat xf|\hat pg\rangle = \langle\hat p\hat
xf|g\rangle$, et $\hat p\hat x = \hat x\hat p - \iu\hbar \ne \hat x\hat p$~: il
n'est pas \omterm{def:b3:quantum-model-systems:hermitian}{hermitien} (sa forme symétrisée $\frac12(\hat x\hat p + \hat p\hat x)$ l'est).
\end{solution}

\begin{solution}{exo:b3:quantum-model-systems:9}
$L = 6 \times 139.7 = \qty{838}{pm}$, $N = 6$~: $\lambda = 8m_ecL^2/(7h) =
\qty{331}{nm}$. L'absorption est dans l'ultraviolet~: l'hexatriène est
incolore.
\end{solution}

\begin{solution}{exo:b3:quantum-model-systems:10}
$\kappa = \sqrt{2m(V_0 - E)}/\hbar$ avec $V_0 - E = \qty{0.20}{eV}$~:
$\kappa_H = \qty{9.82e10}{m^{-1}}$, $\kappa_D = \qty{1.388e11}{m^{-1}}$. Les
préfacteurs se simplifient dans le rapport~: $T_H/T_D = \eu^{2a(\kappa_D - \kappa_H)} =
\eu^{4.06} = 58$, la valeur exacte~: la barrière est épaisse pour les deux
($\kappa a = 4.9$ et 6.9).
\end{solution}

\begin{solution}{exo:b3:quantum-model-systems:11}
Niveaux $E_m = m^2\hbar^2/2m_eR^2$~: $m = 0$ reçoit deux électrons, $m = \pm1$
quatre. Le plus haut occupé est $|m| = 1$, le plus bas vacant $|m| = 2$~: $\Delta E =
3\hbar^2/2m_eR^2 = \qty{9.38e-19}{J}$, $\lambda = \qty{212}{nm}$, dans
l'ultraviolet, là où se trouve l'absorption intense du benzène.
\end{solution}

\begin{solution}{exo:b3:quantum-model-systems:12}
$\langle r\rangle = \frac{4\pi}{\pi a^3}\int_0^\infty r^3\eu^{-2r/a}\dd r =
\frac{4}{a^3}\cdot\frac{3!}{(2/a)^4} = \frac32a$, \qty{79.4}{pm} pour $a = a_0$.
$\langle 1/r\rangle = \frac{4}{a^3}\cdot\frac{1!}{(2/a)^2} = \frac1a$, donc
$\langle V\rangle = -e^2/(4\pi\varepsilon_0a)$. Comme $E_1 =
-e^2/(8\pi\varepsilon_0a)$ (d'après $a = 4\pi\varepsilon_0\hbar^2/\mu e^2$),
$\langle V\rangle = 2E_1$, et l'énergie cinétique moyenne vaut $-E_1$.
\end{solution}

\begin{solution}{pb:b3:quantum-model-systems:1}
\textbf{1.} Chacun des $j$ atomes apporte une orbitale p~; les $j - 1$
carbones et un azote apportent chacun un électron, et l'autre azote son
doublet non liant, puisque la charge du cation est répartie sur la
chaîne~: $N = j + 1$, soit 6,
8 et 10.
\textbf{2.} Six électrons remplissent $n = 1$, 2, 3~: HOMO $n = 3$, LUMO $n = 4$.
\textbf{3.} $\Delta E = E_{N/2+1} - E_{N/2} = (N+1)h^2/8mL^2$ et $\lambda =
hc/\Delta E = 8mcL^2/h(N+1)$.
\textbf{4.} $L_0 = 4l$, $6l$, $8l$~: 559, 838 et \qty{1118}{pm}.
\textbf{5.} 147, 257 et \qty{374}{nm}.
\textbf{6.} Toutes bien trop courtes (d'un facteur 3.6, 2.3 et 1.9)~: la
boîte du modèle est trop petite, puisque $\lambda \propto L^2$~; les
électrons se déplacent sur une distance plus longue que la chaîne
comprise entre les azotes.
\textbf{7.} Avec l'origine au centre, $\psi_n \propto \cos(n\pi u/L)$ pour
$n$ impair et $\sin(n\pi u/L)$ pour $n$ pair ($u = x - L/2$)~: fonctions
paires et impaires de $u$.
\textbf{8.} Si $n + n'$ est pair, $\psi_n$ et $\psi_{n'}$ ont la même parité~;
leur produit multiplié par $u$ est impair et son intégrale sur un
intervalle symétrique est nulle.
\textbf{9.} La HOMO $N/2$ et la LUMO $N/2 + 1$ ont $n + n' = N + 1$,
impair~: transition toujours permise.
\textbf{10.} Deuxième colorant~: HOMO 4, LUMO 5. $4 \to 6$~: somme paire,
interdite. $3 \to 5$~: somme paire, interdite.
\textbf{11.} La suivante permise depuis $n = 4$ est $4 \to 7$, avec un
$\Delta E$ plus grand d'un facteur $(49 - 16)/(25 - 16) = 3.67$~: à $\lambda/3.67$,
\qty{164}{nm} pour le deuxième colorant, loin dans l'ultraviolet.
\textbf{12.} Elle ne dépend que de la symétrie du potentiel par rapport au
centre de la chaîne, que partagent les colorants symétriques réels.
\textbf{13.} $E = hc/\lambda = \qty{3.294e-19}{J} = \qty{2.06}{eV}$.
\textbf{14.} $hc\tilde\omega_e = \qty{0.371}{eV}$.
\textbf{15.} $2hcB = \qty{4.21e-22}{J} = \qty{2.63}{meV}$.
\textbf{16.} $kT = \qty{25.7}{meV}$.
\textbf{17.} Les niveaux de rotation seulement ($2hcB \ll kT$)~; les
excitations de vibration ($15kT$) et électroniques ($80kT$) ne sont
pratiquement pas peuplées.
\textbf{18.} Électroniques~: visible et ultraviolet~; de vibration~:
infrarouge~; de rotation~: micro-ondes (et infrarouge lointain).
\textbf{19.} $L = \sqrt{\lambda h(N+1)/8mc}$ avec $N = 8$~: $L = \qty{1283}{pm}$.
\textbf{20.} $\delta = (1283 - 838)/2 = \qty{222}{pm}$.
\textbf{21.} $L = 559 + 445 = \qty{1003}{pm}$, $\lambda = \qty{474}{nm}$, 9.5\,\%
au-dessous de 524 nm.
\textbf{22.} $L = 1118 + 445 = \qty{1562}{pm}$, $\lambda = \qty{732}{nm}$, 2.9\,\%
au-dessus de 711 nm.
\textbf{23.} Dans les cycles aromatiques qui portent les azotes~: le
système $\pi$ ne s'arrête pas aux atomes d'azote, et la boîte effective
englobe une partie de chaque cycle.
\textbf{24.} $\delta/l = 222/139.7 = 1.6$~: \textbf{la boîte se prolonge
d'environ \qty{222}{pm}, une liaison et demie, au-delà de chaque azote},
et avec cette seule longueur le modèle reproduit les trois couleurs à
10\,\% près.
\end{solution}
